% Related lecture: SC4002-L06
% Source: Week 6-Constitency Parsing, PDF/slide pp. 8, 15--18, 21--31
% E01--E03: lecture examples restated with solutions.
% Human verification: learner requested notes on 2026-09-16 after the explanation.
% E01--E03 were explained and source-checked by assistant, not independently solved.

\exercisegroup{SC4002-L06-EX}{第 06 讲：Constituency Grammars and Parsing}

\exercisemeta
  {SC4002-L06-EX}
  {2026-09-16}
  {Week 6-Constitency Parsing pp.~8, 15--18, 21--31；本地课堂字幕 61156}
  {E01--E03 为讲义例题重述与已核对解答，非 Tutorial}
  {用户听完讲解后请求整理；例题由助手讲解及核对}

\exercisequestion{SC4002-L06-E01}{CFG 例题：I prefer a morning flight}

\textbf{类型与来源：}lecture worked example，讲义 p.~8；按该页产生式与树重述。

\textbf{对应知识点：}\relatedtopic{SC4002-L06-T01}{Constituency、CFG 与 Treebanks}。

\subsubsection*{题意与条件}

利用以下规则解析 \texttt{I prefer a morning flight}，说明 NP 与 Nominal 的区别：
\[
\begin{aligned}
  S&\rightarrow NP\ VP,& NP&\rightarrow Pronoun\mid Det\ Nominal,\\
  VP&\rightarrow Verb\ NP,& Nominal&\rightarrow Nominal\ Noun\mid Noun,\\
  Pronoun&\rightarrow\texttt{I},& Verb&\rightarrow\texttt{prefer},\\
  Det&\rightarrow\texttt{a},& Noun&\rightarrow\texttt{morning}\mid\texttt{flight}.
\end{aligned}
\]

\begin{workedsolution}
主干 $S\rightarrow NP\ VP$ 分出主语 \texttt{I} 和 VP \texttt{prefer a morning flight}。VP 再由 Verb 与宾语 NP 构成；宾语 NP 使用 $NP\rightarrow Det\ Nominal$，其中 Nominal 覆盖 \texttt{morning flight}。最后以 $Nominal\rightarrow Nominal\ Noun$ 展开，左侧 Nominal 再经 Noun 产生 \texttt{morning}，右侧 Noun 产生 \texttt{flight}。

\begin{center}
  % Original TikZ redraw from SC4002 Week 6, slide 8.
% Follows the displayed tree and productions (Nominal -> Nominal Noun).
\begin{tikzpicture}[x=1cm,y=0.9cm,
  every node/.style={font=\small,inner sep=2pt},
  edge/.style={draw=noteBlue,thick}]
  \node (s) at (3.5,0) {S};
  \node (np1) at (0.7,-0.9) {NP};
  \node (vp) at (5.1,-0.9) {VP};
  \node (pro) at (0.7,-1.8) {Pronoun};
  \node (v) at (3.1,-1.8) {Verb};
  \node (np2) at (6.4,-1.8) {NP};
  \node (det) at (4.9,-2.7) {Det};
  \node (nom) at (7.5,-2.7) {Nominal};
  \node (nom1) at (6.5,-3.6) {Nominal};
  \node (n2) at (8.8,-3.6) {Noun};
  \node (n1) at (6.5,-4.5) {Noun};
  \node[text=ntuRed] (i) at (0.7,-5.3) {I};
  \node[text=ntuRed] (prefer) at (3.1,-5.3) {prefer};
  \node[text=ntuRed] (a) at (4.9,-5.3) {a};
  \node[text=ntuRed] (morning) at (6.5,-5.3) {morning};
  \node[text=ntuRed] (flight) at (8.8,-5.3) {flight};
  \foreach \parent/\child in {s/np1,s/vp,np1/pro,pro/i,vp/v,v/prefer,vp/np2,np2/det,det/a,np2/nom,nom/nom1,nom/n2,nom1/n1,n1/morning,n2/flight}
    \draw[edge] (\parent) -- (\child);
\end{tikzpicture}


  {\small 原创重绘；依据 SC4002 Week 6 p.~8 的产生式与树。}
\end{center}

树叶依次还原原句。Nominal 覆盖 \texttt{morning flight}，加入 \texttt{a} 后经给定规则成为 NP；类别与覆盖范围都需按树区分。
\end{workedsolution}

\exercisequestion{SC4002-L06-E02}{结构歧义：PP Attachment 与 Coordination}

\textbf{类型与来源：}lecture examples，讲义 pp.~15--18；括号为对原图结构的简化表达，不展开词汇层。

\textbf{对应知识点：}\relatedtopic{SC4002-L06-T02}{Structural Ambiguity}。

\subsubsection*{题意}

说明 \texttt{I shot an elephant in my pajamas} 中 PP 的两种附着方式，以及 \texttt{old men and women} 中形容词的两种作用范围。再解释为什么 CFG 合法不等于语义合理。

\begin{workedsolution}
\begin{enumerate}
  \item \texttt{I shot [an elephant [in my pajamas]]}：PP 在宾语 NP 的名词性部分内部，表示大象穿着睡衣。\texttt{I [[shot an elephant] [in my pajamas]]}：PP 接到 VP 层，通常表示说话者穿着睡衣进行射击。
  \item \texttt{[old men] and women} 只说明 men 年老，women 的年龄未说明；\texttt{old [men and women]} 同时修饰两组人。
\end{enumerate}
同理，\texttt{A policeman shot the criminal with a revolver} 的 PP 可描述射击工具，也可描述罪犯携带的物品。\texttt{We saw the Eiffel Tower flying to Paris} 通常希望表达 ``我们飞往巴黎时看到铁塔''，而错误附着会产生 ``铁塔在飞'' 的解释。语法规则约束组合形式，语义、常识与语境才帮助在候选树中消歧；不应把不常见的语义直接等同于文法不允许。
\end{workedsolution}

\exercisequestion{SC4002-L06-E03}{CKY：Book the flight through Houston}

\textbf{类型与来源：}lecture worked example，讲义 pp.~21--31；三条 S 推导按 p.~30 的 $S_1,S_2,S_3$ 编号。以下公式与简化树是对该例的展开，不是额外编造的句子。

\textbf{对应知识点：}\relatedtopic{SC4002-L06-T03}{CNF 与 CKY}。

\subsubsection*{题意与文法}

用 CNF grammar 解析 \texttt{Book the flight through Houston}，填出 chart，解释最后一列，并回溯整句的三种 S 推导。边界为
\[
  0\ \texttt{Book}\ 1\ \texttt{the}\ 2\ \texttt{flight}\ 3\
  \texttt{through}\ 4\ \texttt{Houston}\ 5.
\]
讲义 p.~21 的二元规则如下；词汇规则只列本句所需者，其余词不会参与本次推导。句首 \texttt{Book} 按讲义匹配词汇项 \texttt{book}。
\[
\begin{aligned}
  S&\rightarrow NP\ VP\mid X1\ VP\mid Verb\ NP
       \mid X2\ PP\mid Verb\ PP\mid VP\ PP,\\
  X1&\rightarrow Aux\ NP,\qquad X2\rightarrow Verb\ NP,\\
  NP&\rightarrow Det\ Nominal,\\
  Nominal&\rightarrow Nominal\ Noun\mid Nominal\ PP,\\
  VP&\rightarrow Verb\ NP\mid X2\ PP\mid Verb\ PP\mid VP\ PP,\\
  PP&\rightarrow Preposition\ NP.
\end{aligned}
\]
\begin{center}
\small
\begin{tabularx}{\textwidth}{@{}lX@{}}
  \toprule
  词汇项 & 可直接产生它的非终结符 \\
  \midrule
  \texttt{book} & S, VP, Verb, Nominal, Noun \\
  \texttt{the} & Det \\
  \texttt{flight} & Nominal, Noun \\
  \texttt{through} & Preposition \\
  \texttt{Houston} & NP, Proper-Noun \\
  \bottomrule
\end{tabularx}
\end{center}
S、VP、Nominal 的直接词汇规则来自消去单链，例如 $S\Rightarrow VP\Rightarrow Verb\Rightarrow\texttt{book}$；因此表内类别不只是 POS tags。

\begin{workedsolution}
\textbf{短区间。}先填五个单词对应的 $C[i,i+1]$。随后
\[
\begin{aligned}
  Det[1,2]+Nominal[2,3]&\xRightarrow{NP\rightarrow Det\ Nominal}NP[1,3],\\
  Verb[0,1]+NP[1,3]&\Longrightarrow S[0,3],\ VP[0,3],\ X2[0,3].
\end{aligned}
\]
第二行同时命中三条规则，不能只保留其中一个父类别。完整 chart 的\emph{类别集合}如下；重复的同类推导另存 backpointers，不通过重复类别名表达。

\begin{center}
  % Original TikZ redraw; SC4002 Week 6 slides 23--30, verified with slide 21 rules.
% Category sets only; three distinct S backpointers are explained in E03.
\begin{tikzpicture}[x=1cm,y=1cm,
  every node/.style={font=\small,align=center,inner sep=2pt}]
  \foreach \i in {0,...,4}{
    \pgfmathtruncatemacro{\firstj}{\i+1}
    \foreach \j in {\firstj,...,5}{
      \draw[draw=noteBlue!60,fill=noteBlue!5]
        ({2.45*(\j-1)},{-1.32*\i}) rectangle ++(2.45,-1.32);
      \node[font=\scriptsize,text=black!65]
        at ({2.45*(\j-0.5)},{-1.32*\i-1.12}) {$[\i,\j]$};
    }
  }
  \foreach \j/\word in {1/Book,2/the,3/flight,4/through,5/Houston}
    \node at ({2.45*(\j-0.5)},0.38) {\word};
  \node at (1.225,-0.46) {S, VP, Verb\\Nominal, Noun};
  \node at (3.675,-0.46) {$\varnothing$};
  \node at (6.125,-0.46) {S, VP, X2};
  \node at (8.575,-0.46) {$\varnothing$};
  \node[text=ntuRed] at (11.025,-0.46) {S, VP, X2\\\scriptsize S: 3 derivations};
  \node at (3.675,-1.78) {Det};
  \node at (6.125,-1.78) {NP};
  \node at (8.575,-1.78) {$\varnothing$};
  \node at (11.025,-1.78) {NP};
  \node at (6.125,-3.10) {Nominal, Noun};
  \node at (8.575,-3.10) {$\varnothing$};
  \node at (11.025,-3.10) {Nominal};
  \node at (8.575,-4.42) {Preposition};
  \node at (11.025,-4.42) {PP};
  \node at (11.025,-5.74) {NP\\Proper-Noun};
  \node[anchor=north west,align=left,text width=4.3cm,font=\small]
    at (0,-2.85) {边界编号：$0,1,\ldots,5$。\\
    $[1,3]$: \texttt{the flight}\\
    $[3,5]$: \texttt{through Houston}\\[0.6em]
    从左向右逐列填表；\\每列从下向上。};
\end{tikzpicture}


  {\small 原创重绘；依据讲义 pp.~21--30。$\varnothing$ 表示当前文法没有合法成分。}
\end{center}

\textbf{最后一列。}除已填好的 $C[4,5]=\{NP,Proper\text{-}Noun\}$ 外，由下向上计算：
\begin{enumerate}
  \item $C[3,5]$：以 $k=4$ 切为 \texttt{through} 与 \texttt{Houston}，由 $PP\rightarrow Preposition\ NP$ 得到 PP。
  \item $C[2,5]$：以 $k=3$ 切为 \texttt{flight} 与 \texttt{through Houston}，由 $Nominal\rightarrow Nominal\ PP$ 得到 Nominal；$k=4$ 的左格为空。
  \item $C[1,5]$：必须检查三个切分点，只有 $k=2$ 成功：
\end{enumerate}

\begin{center}
\small
\begin{tabularx}{\textwidth}{@{}cllX@{}}
  \toprule
  $k$ & 左格 & 右格 & 是否能合并 \\
  \midrule
  2 & $C[1,2]=\{Det\}$ & $C[2,5]=\{Nominal\}$ & 有 $NP\rightarrow Det\ Nominal$，得到 NP \\
  3 & $C[1,3]=\{NP\}$ & $C[3,5]=\{PP\}$ & 无 $NP\rightarrow NP\ PP$，也无其他匹配规则 \\
  4 & $C[1,4]=\varnothing$ & $C[4,5]$ 非空 & 左格为空，不成立 \\
  \bottomrule
\end{tabularx}
\end{center}

所以 $C[1,5]=\{NP\}$。这里特别说明：$k=3$ 的两格都非空仍不能合并，因为不得临时添加文法中没有的规则。

\textbf{整句的三条 S 推导。}在 $C[0,5]$ 中，$k=1$ 命中一条 S 规则，$k=3$ 命中两条；$k=2,4$ 均因左格为空失败：
\[
\begin{aligned}
  S_1[0,5]&\rightarrow Verb[0,1]\ NP[1,5],\\
  S_2[0,5]&\rightarrow VP[0,3]\ PP[3,5],\\
  S_3[0,5]&\rightarrow X2[0,3]\ PP[3,5].
\end{aligned}
\]
下标仅区分 S 的三个来源，不是三个新文法类别。整个格子的类别集合为 $\{S,VP,X2\}$，但只有以起始符号 S 为根的推导是完整句子解析；VP 本身也有多个推导。

\textbf{恢复原文法。}每个 S 都恢复 $S\rightarrow VP$。第三种还展开并删除辅助 $X2$；NP、Nominal 等处被压缩的词汇单链也要恢复。省略词汇层后，三棵树的主要分组为：
\begin{enumerate}
  \item $S_1$：$[_{VP}\ Verb\ [_{NP}\ Det\ [_{Nominal}\ Nominal\ PP]]]$。PP 在宾语的名词性部分内，修饰 flight。
  \item $S_2$：$[_{VP}\ [_{VP}\ Verb\ NP]\ PP]$。先构成 \texttt{Book the flight} 这个 VP，再附加 PP。
  \item $S_3$：$[_{VP}\ Verb\ NP\ PP]$。来自原文法的三分支规则 $VP\rightarrow Verb\ NP\ PP$。
\end{enumerate}
第二、三种都将 PP 置于 VP 层面，但一个包含嵌套 VP，一个直接有三个子节点。三棵句法树不必对应三种不同语义。仅存集合 $\{S,VP,X2\}$ 不足以恢复这些差异，必须保留每个成功组合的 backpointer。
\end{workedsolution}
